% Part: turing-machines
% Chapter: machines-computations
% Section: variants

\documentclass[../../../include/open-logic-section]{subfiles}

\begin{document}

\olfileid{tur}{mac}{var}
\olsection{Varian Mesin Turing}

Sejatine ana akeh cara kanggo netepake mesin Turing,
lan definisi kita mung salah sijine. Ing sawetara bab,
definisi kita luwih luwes tinimbang definisi liyane.
Kita ngidinake alfabet winates apa wae, dene definisi
kang luwih winates bisa uga mung ngidinake rong simbol
pita, $\TMstroke$ lan~$\TMblank$. Kita ngidinake mesin
nulis simbol ing pita lan obah bebarengan, dene
definisi liyane mung ngidinake nulis utawa obah.
Kita ngidinake nulis tanpa mindhah sirah maca-nulis,
dene definisi liyane ora nganggo ``instruksi''
$\TMstay$. Ing bab liyane, definisi kita luwih
winates. Kita nganggep pita tanpa wates mung ing siji
arah, dene definisi liyane ngidinake pita tanpa wates
ing sisih kiwa lan tengen. Malah bisa uga diidinake
pita kapisah kanthi cacah winates pira wae, utawa kothak-kothak
kang mbentuk kisi tanpa wates. Kita nggambarake
himpunan instruksi mesin Turing nganggo fungsi transisi;
definisi liyane nganggo relasi transisi, saengga mesin
bisa nduwé luwih saka siji instruksi kang bisa ditindakake
ing sawijining kahanan.

Panglonggaran definisi kang pungkasan iki narik kawigaten.
Ing definisi kita, nalika mesin ana ing kahanan~$q$
lan maca simbol~$\sigma$, $\delta(q, \sigma)$
nemtokake simbol anyar, kahanan anyar, lan panggonan
anyar sirah maca-nulis. Nanging yen himpunan instruksi
diidinake dadi relasi antarane pasangan kahanan-simbol
saiki $\tuple{q, \sigma}$ lan tripel kahanan-simbol-arah
anyar $\tuple{q', \sigma', D}$, tumindake mesin Turing
bisa uga ora ditemtokake kanthi tunggal---relasi
instruksi bisa ngemot $\tuple{q, \sigma, q', \sigma', D}$
lan $\tuple{q, \sigma, q'', \sigma'', D'}$ bebarengan.
Ing kahanan iki kita nduwé mesin Turing
\emph{nondeterministik}. Mesin kaya mangkono nduwé
peran wigati ing teori kompleksitas komputasi.

Uga ana pranatan kang beda-beda bab kapan mesin Turing
mandheg: miturut kita, mesin mandheg yen fungsi transisine
ora ditegesi; definisi liyane mbutuhake mesin mlebu
kahanan mandheg mligi. Ing \olref[hal]{sec} kita wis
njlentrehake apa sebabe syarat kahanan mandheg mligi
ora ngowahi apa kang bisa diitung mesin Turing.
Amarga pita mesin Turing kita tanpa wates mung ing siji
arah, ana kahanan nalika mesin ora bisa nindakake
instruksi kanthi lumrah: nalika maca kothak paling kiwa
nanging didhawuhi obah menyang kiwa. Miturut definisi
kita, sirah mung tetep ing panggonane tinimbang
``tiba saka pita'', nanging kita uga bisa netepake
yen mesin mandheg nalika iku kedadeyan. Definisi iki
uga padha daya komputasine: tumindake mesin Turing kang
mandheg nalika nyoba obah kiwa saka kothak~$0$ bisa
kita tiru kanthi mbusak saben transisi
$\delta(q,\TMendtape) =
\tuple{q',\sigma,\TMleft}$---mula tinimbang nyoba
obah kiwa nalika maca~$\TMendtape$, mesin mandheg.\footnote{Cara iki
ora \emph{sakabehe} lumaku, amarga kita ora dilarang
nulis lan maca $\TMendtape$ ing kothak saliyane
kothak~$0$ (delengen \olref[una]{ex:mover}). Iki bisa
diatasi kanthi nambah simbol~$\TMendtape'$ kapindho
kanggo digunakake ing kahanan kaya mangkono. Cathetan cakupan SOL6-F053: cara pambusakan instruksi iki uga nganggep panandha pucuk asli ing kothak nol tetep lestari. Yen panandha iku bisa ditimpa, dhisik tiru mesin kanthi pangodean kang ngenali wates kiwa asli kanthi mligi.}

Uga ana cara liya kanggo nggambarake wilangan
(lan kanthi mangkono fungsi input-output kang diitung
mesin Turing): kita nganggo representasi kanthi siji
simbol, nanging representasi biner uga bisa digunakake.
Cara biner iki mbutuhake rong simbol saliyane
$\TMblank$ lan~$\TMendtape$.

Ana kasunyatan kang narik kawigaten: varian-varian iki
ora ngowahi himpunan fungsi kang bisa diitung dening
mesin Turing. \emph{Yen sawijining fungsi bisa diitung
dening mesin Turing miturut siji definisi, fungsi iku
uga bisa diitung miturut kabeh definisi mau.} Cathetan cakupan SOL6-F054: kanggo fungsi nondeterministik, kabeh cabang kang nampa lan menehi output kanggo input padha kudu menehi nilai padha, lan yen nilai ditegesi kudu ana cabang kang nampa lan mandheg. Cabang liyane bisa nolak utawa ora mandheg. Penelusuran adil miturut dawa langkah bisa nemokake cabang panriman; relasi kang menehi maneka output kanggo input padha ora otomatis fungsi.

Kita ora bakal mbahas bukti rinci bab iki. Iki siji
tuladhane: pita kang tanpa wates ing rong arah utawa
pita luwih saka siji ora nambah daya komputasi.
Alasane, kanthi garis gedhe, mesin Turing kang nduwé
siji pita tanpa wates ing siji arah bisa niru pita
luwih saka siji utawa pita tanpa wates ing rong arah.
Contone, kanthi nambah kahanan lan instruksi, kita
bisa ``nerjemahake'' program mesin kang nduwé pirang-pirang
pita utawa pita tanpa wates ing rong arah dadi program
mesin kang mung nduwé siji pita tanpa wates ing siji
arah. Mesin kang diasilake bisa nggunakake kothak
genep kanggo kothak-kothak pita~$1$ (utawa kothak-kothak
``positif'' saka pita tanpa wates rong arah), lan
kothak ganjil kanggo kothak-kothak pita~$2$ (utawa
kothak-kothak ``negatif'').

\end{document}
